Resuelve las siguientes ecuaciones diferenciales o problemas de valor inicial mediante un cambio de variable. Explica por qué el cambio elegido simplifica la ecuación y comprueba la solución en las variables originales. Trabaja con soluciones reales de clase $C^1$ en intervalos abiertos; indica las restricciones de dominio y las ramas que se pierdan al dividir. En los problemas de valor inicial, da la solución en su intervalo máximo dentro del dominio indicado. Se admiten respuestas implícitas cuando la inversión no sea elemental.
\begin{enumerate}[label=\alph*)]
\item $yy'=x(1-y^2)$, \quad $y(0)=2$.
\item $xy'+y\ln y=y\ln x$, \quad $x>0$, $y>0$.
\item $y'=1+e^{-y}\cos x$.
\item $(1+x^2)y'=2xy+y^2$.
\item $(1+3y^2)y'+\dfrac{2}{x}(y+y^3)=x$, \quad $x>0$.
\item $y'=\bigl(1+(y+2x)^2\bigr)\arctan(y+2x)-2$, \quad $y(0)=1$.
\item $2y'=1+\sin(2y-x)$.
\item $y'=1+\dfrac{(y-x)^2-1}{y-x+2}$.
\item $y'=\dfrac yx+\dfrac{x(y+x)}{x^2+y^2}$, \quad $x\ne0$.
\item $y'=\dfrac{y+1}{x-2}+\dfrac{(x-2)^2}{(x-2)^2+(y+1)^2}$, \quad $x\ne2$.
\item $xy'-2y=x^3e^{-y/x^2}$, \quad $x>0$.
\item $y'=y+e^x\bigl(1+e^{-2x}y^2\bigr)^{3/2}$.
\item $\sin(2y)y'=\sin(2x)$, \quad $y(0)=0$. Determina todas las soluciones $C^1$ del PVI y justifica qué ocurre cuando $\sin(2y)=0$.
\item $\left(\dfrac1y+x\right)y'+y=\dfrac{1}{1+x^2}$, \quad $x>0$, $y>0$, \quad $y(1)=1$.
\item $y'=\dfrac{1}{x-y^2+2y}$. Puede resultar útil considerar $x$ como función de $y$.
\end{enumerate}
